\documentclass[12pt,a4paper]{article}
\usepackage[utf8]{inputenc}
\usepackage[spanish,es-noquoting]{babel}
\usepackage[T1]{fontenc}
\usepackage{lmodern}
\usepackage{geometry}
\geometry{top=2cm, bottom=2cm, left=2.5cm, right=2.5cm}
\usepackage{xcolor}
\usepackage{listings}
\usepackage{enumitem}
\usepackage{fancyhdr}
\usepackage{hyperref}
\usepackage{tikz}
\usetikzlibrary{arrows.meta, positioning, calc, shapes.geometric}

\definecolor{codegreen}{rgb}{0,0.6,0}
\definecolor{backcolour}{rgb}{0.95,0.95,0.92}

\lstset{
    language=C,
    backgroundcolor=\color{backcolour},
    commentstyle=\color{codegreen},
    keywordstyle=\color{blue},
    basicstyle=\ttfamily\small,
    breaklines=true,
    frame=single,
    numbers=left,
    tabsize=4,
    extendedchars=true,
    showstringspaces=false,
    literate={á}{{\'a}}1 {é}{{\'e}}1 {í}{{\'i}}1 {ó}{{\'o}}1 {ú}{{\'u}}1 {ñ}{{\~n}}1
}

\pagestyle{fancy}
\fancyhf{}
\lhead{Monitorías Sistemas Operativos 2026-I}
\rhead{Capítulo 1: Jerarquía de Procesos}
\cfoot{\thepage}

\title{\textbf{Taller de Lógica: Jerarquía y Creación de Procesos}}
\author{Malak Sanchez \\ \small Monitorías de Sistemas Operativos}
\date{\today}

\begin{document}

\maketitle

\section*{Introducción}
Este taller se enfoca puramente en la lógica de creación y herencia de procesos mediante la llamada al sistema \texttt{fork()}. No se requiere sincronización ni comunicación avanzada; el objetivo es que el estudiante domine la visualización de estructuras complejas y el comportamiento de los procesos en bifurcación.

\section*{Bloque 1: De Diagramas a Código}
\textit{Desarrolle el código necesario para replicar exactamente las siguientes estructuras de jerarquía.}

\subsection*{Ejercicio 1}
Un proceso padre que genera exactamente 7 procesos hijos de forma paralela.
\begin{center}
\begin{tikzpicture}[node distance=1cm, every node/.style={draw, circle, inner sep=2.2pt, fill=black}]
    \node (P) {};
    \foreach \i in {1,...,7} {
        \node (H\i) [below=of P, xshift=\i*1cm - 4cm] {};
        \draw (P) -- (H\i);
    }
\end{tikzpicture}
\end{center}

\subsection*{Ejercicio 2}
Un proceso raíz con 3 hijos: el primero tiene 1 hijo, el segundo tiene 2 hijos y el tercero tiene 3 hijos.
\begin{center}
\begin{tikzpicture}[node distance=1.2cm, every node/.style={draw, circle, inner sep=2.2pt, fill=black}]
    \node (R) {};
    
    % Level 1
    \node (H1) [below left=of R, xshift=-1cm] {};
    \node (H2) [below=of R] {};
    \node (H3) [below right=of R, xshift=1cm] {};
    \draw (R) -- (H1); \draw (R) -- (H2); \draw (R) -- (H3);
    
    % Level 2
    \node (N1) [below=of H1] {}; \draw (H1) -- (N1);
    
    \node (N2a) [below left=of H2, xshift=0.8cm] {}; 
    \node (N2b) [below right=of H2, xshift=-0.8cm] {};
    \draw (H2) -- (N2a); \draw (H2) -- (N2b);
    
    \node (N3a) [below left=of H3, xshift=0.8cm] {};
    \node (N3b) [below=of H3] {};
    \node (N3c) [below right=of H3, xshift=-0.8cm] {};
    \draw (H3) -- (N3a); \draw (H3) -- (N3b); \draw (H3) -- (N3c);
\end{tikzpicture}
\end{center}

\newpage

\subsection*{Ejercicio 3}
Replique la jerarquía de la imagen: Un proceso raíz con 5 hijos, donde el tercer hijo crea a su vez 4 hijos. Los extremos de esta última generación (el primero y el cuarto) crean un proceso final cada uno.
\begin{center}
\begin{tikzpicture}[node distance=1cm, every node/.style={draw, circle, inner sep=2pt, fill=black}, scale=0.8]
    \node (R) {};
    \foreach \i in {1,...,5} {
        \node (H\i) [below=of R, xshift=\i*1.5cm - 4.5cm] {};
        \draw (R) -- (H\i);
    }
    \foreach \j in {1,...,4} {
        \node (N\j) [below=of H3, xshift=\j*1cm - 2.5cm] {};
        \draw (H3) -- (N\j);
    }
    \node (B1) [below=of N1] {}; \draw (N1) -- (B1);
    \node (B4) [below=of N4] {}; \draw (N4) -- (B4);
\end{tikzpicture}
\end{center}

\subsection*{Ejercicio 4}
Un padre con 2 hijos. Cada hijo crea 1 nieto, y cada nieto crea 1 bisnieto.
\begin{center}
\begin{tikzpicture}[node distance=0.8cm, every node/.style={draw, circle, inner sep=2pt, fill=black}]
    \node (R) {};
    \node (H1) [below left=of R] {}; \node (H2) [below right=of R] {};
    \draw (R) -- (H1); \draw (R) -- (H2);
    \node (N1) [below=of H1] {}; \node (N2) [below=of H2] {};
    \draw (H1) -- (N1); \draw (H2) -- (N2);
    \node (B1) [below=of N1] {}; \node (B2) [below=of N2] {};
    \draw (N1) -- (B1); \draw (N2) -- (B2);
\end{tikzpicture}
\end{center}

\subsection*{Ejercicio 5}
Desarrolle una función recursiva \texttt{void expandir(int n)} que genere un árbol binario perfecto de profundidad 3. Dibuje el árbol resultante e indique el número total de procesos.

\newpage

\section*{Bloque 2: De Código a Diagramas}
\textit{Analice los siguientes fragmentos de código y dibuje el árbol de procesos resultante. Indique cuántos procesos totales se crearon (incluyendo al padre original).}

\subsection*{Ejercicio 6}
\begin{lstlisting}
int n = 2;
for(int i=0; i < n; i++) {
    fork();
    fork();
}
\end{lstlisting}

\subsection*{Ejercicio 7}
\begin{lstlisting}
if(fork() || fork()) {
    fork();
}
\end{lstlisting}

\subsection*{Ejercicio 8}
\begin{lstlisting}
if(fork() && fork()) {
    fork();
}
\end{lstlisting}

\subsection*{Ejercicio 9}
\begin{lstlisting}
pid_t p;
p = fork();
if(p != 0) {
    fork();
}
fork();
\end{lstlisting}

\newpage

\subsection*{Ejercicio 10}
Analice qué sucede en este código y dibuje la jerarquía generada.
\begin{lstlisting}
void generar(int n) {
    if(n <= 0) return;
    if(fork() == 0) {
        generar(n-1);
    }
    fork();
}

int main() {
    generar(2);
    return 0;
}
\end{lstlisting}

\subsection*{Ejercicio 11}
\begin{lstlisting}
pid_t p;
p = fork();
if(p = 0) {
    fork();
}
\end{lstlisting}

\end{document}
